Rozdział ten jest poświęcony szczegółowej analizie uzyskanych rezultatów eksperymentów w odniesieniu do postawionych pytań badawczych oraz hipotez. Wskazano w nim powiązania pomiędzy architekturą silników a ich charakterystyką wydajnościową, a także sformułowano praktyczne wnioski inżynierskie.

\section{Wpływ architektury bazy na wyniki}
Wykazano bezpośrednią korelację pomiędzy wewnętrzną architekturą silników bazodanowych a ich charakterystyką opóźnień. W szczególności poddano analizie zjawisko tzw. burz kompaktowania (ang. \emph{compaction storms}) występujących w silnikach opartych na strukturze drzewa LSM (LevelDB oraz RocksDB) \cite{oneil1996}:
\begin{itemize}
    \item Podczas intensywnego zapisu dane są początkowo buforowane w strukturach pamięciowych (MemTable), a następnie zrzucane na dysk w postaci plików SST.
    \item W tle uruchamiane są asynchroniczne procesy kompaktowania scalające pliki na kolejnych poziomach drzewa LSM. Nakładanie się tych operacji I/O w tle z bieżącymi zapytaniami klienckimi powoduje chwilowe wysycenie przepustowości dyskowej.
    \item Skutkuje to znaczącym wzrostem opóźnień ogonowych (ang. \emph{tail latency}), widocznym w postaci nagłych pików w centylach opóźnień p95 oraz p99, przy jednoczesnym zachowaniu niskiego opóźnienia średniego (p50). Zjawisko to zestawiono z pracą baz in-memory (Redis, Memcached, Tarantool), które wykazują stabilną charakterystykę opóźnień niezależnie od czasu trwania testu.
\end{itemize}

\section{Wpływ środowiska uruchomieniowego}

Porównanie wydajności w środowisku natywnym Linux oraz wirtualizacji WSL 2 wykazało istotne różnice w zależności od profilu obciążenia (I/O-bound vs CPU-bound):
\begin{itemize}
    \item \textbf{Wydajność podsystemu sieciowego (Loopback):} Nowoczesna implementacja wirtualnego loopbacku w WSL 2, wykorzystująca mechanizm AF\_VSOCK do bezpośredniej komunikacji między jądrem Linux a systemem operacyjnym Windows, pozwoliła na uzyskanie bardzo niskich narzutów. Dla systemów in-memory operujących w sieci (Redis, Memcached) spadek przepustowości nie przekroczył 8\%, co potwierdza wysoką dojrzałość WSL 2 w zastosowaniach deweloperskich.
    \item \textbf{Wydajność podsystemu dyskowego (VHDX):} Wirtualny dysk WSL 2 zapisywany jest jako plik w formacie VHDX na systemie plików NTFS systemu Windows. Każde wywołanie funkcji synchronizacji dyskowej (\texttt{fsync}) przez systemy plików wewnątrz Linuksa wymaga translacji na wywołania systemowe NTFS i fizycznego opróżnienia buforów dysku SSD hosta. W efekcie bazy transakcyjne z silnym zapisem WAL (FoundationDB, Couchbase) zanotowały w WSL 2 spadek przepustowości zapisu o 25\%--40\% w odniesieniu do natywnego systemu Ubuntu.
    \item \textbf{Jitter szeregowania procesów (CPU Scheduling):} Dodatkowa warstwa wirtualizacji (mikro-hiperwizor Hyper-V) wprowadza narzut na przełączanie kontekstu procesora. Spowodowało to podwyższenie opóźnień ogonowych p95 i p99 o 12\%--18\% dla wszystkich systemów w WSL 2. Produkcyjne wdrożenia baz danych wymagają zatem natywnego środowiska Linux w celu zagwarantowania stabilności czasu odpowiedzi.
\end{itemize}

\section{Wnioski praktyczne -- rekomendacje wdrożeniowe}

Na podstawie przeprowadzonych badań sformułowano następujące zalecenia architektoniczne:
\begin{itemize}
    \item \textbf{Szybki cache in-memory bez trwałości dyskowej:} Rekomendowane jest użycie systemu \textbf{Memcached}. Charakteryzuje się on najniższymi opóźnieniami i najwyższą wydajnością dzięki wielowątkowej obsłudze żądań.
    \item \textbf{Baza Key-Value z obsługą struktur danych i trwałością:} Standardem pozostaje system \textbf{Redis}, łączący bogatą funkcjonalność z bardzo stabilnymi czasami odpowiedzi. W przypadku projektów wymagających dedykowanego serwera aplikacji wewnątrz bazy, doskonałym wyborem jest \textbf{Tarantool}.
    \item \textbf{Lokalny zapis bez narzutu sieciowego w aplikacjach jednowątkowych:} Należy wybrać bibliotekę \textbf{RocksDB} zamiast starszej LevelDB. Oferuje ona znacznie lepsze mechanizmy optymalizacji kompaktowania LSM oraz mechanizm Column Families.
    \item \textbf{Lokalna baza danych z dostępem wieloprocesowym:} Jedynym w pełni sprawnym rozwiązaniem wbudowanym jest \textbf{BerkeleyDB}, która poprzez pamięć dzieloną obsługuje współbieżny dostęp transakcyjny z wielu niezależnych procesów systemowych.
    \item \textbf{Duże systemy rozproszone o dynamicznej strukturze dokumentów:} Zalecany jest system \textbf{Couchbase}, oferujący natywny sharding, replikację oraz elastyczność formatu JSON.
    \item \textbf{Systemy krytyczne wymagające globalnej spójności transakcyjnej (ACID):} Rekomendowanym wyborem jest \textbf{FoundationDB}, gwarantująca spójność transakcyjną na poziomie całego klastra rozproszonego.
\end{itemize}

\section{Ograniczenia badań}
Prezentowane wyniki badań, mimo zachowania rygoru metodologicznego, obarczone są pewnymi ograniczeniami wynikającymi z przyjętej koncepcji eksperymentu oraz warunków sprzętowych:
\begin{itemize}
    \item \textbf{Kolokacja procesów klienckich i serwera:} Uruchomienie aplikacji testowej (klienta) oraz silników baz danych (serwera) na tej samej fizycznej maszynie sprzętowej, mimo zastosowania limitów cgroups (Docker), generuje wzajemną rywalizację o zasoby pamięci podręcznej procesora (CPU cache L3) oraz przepustowość magistrali pamięci RAM. W rzeczywistych wdrożeniach produkcyjnych aplikacja i baza danych są zazwyczaj rozdzielone fizycznie i połączone dedykowaną siecią LAN.
    \item \textbf{Skalowanie pionowe zamiast poziomego:} Eksperymenty koncentrowały się na badaniu wydajności w obrębie jednego węzła (skalowanie pionowe -- pionowa alokacja rdzeni CPU). Ogranicza to możliwość bezpośredniego przełożenia uzyskanych wyników na zachowanie baz w klastrach wielowęzłowych (skalowanie poziome), gdzie kluczowym czynnikiem limitującym staje się opóźnienie sieci WAN/LAN, narzut protokołów spójności (np. Raft, Paxos) czy proces partycjonowania danych.
    \item \textbf{Syntetyczny profil obciążenia:} Zastosowany benchmark opiera się na generowaniu powtarzalnych, jednolicie rozłożonych losowych operacji o stałym rozmiarze rekordu (128 bajtów oraz 4 KB). W rzeczywistych systemach produkcyjnych rozkład zapytań jest zazwyczaj skośny (np. zgodny z rozkładem Zipfa, zasada Pareto), a rozmiary przechowywanych obiektów charakteryzują się dużą wariancją, co bezpośrednio wpływa na efektywność buforowania i dynamikę kompaktowania danych na dysku.
    \item \textbf{Uproszczenia emulacji struktur danych:} Operacje na kolejce FIFO oraz dokumentach JSON dla systemów nieposiadających natywnego wsparcia były emulowane po stronie aplikacji w języku Python. Narzut serializacji i deserializacji obiektów w jednowątkowym interpreterze CPython mógł w niektórych konfiguracjach stanowić dominującą składową czasu wykonania zapytania, co uwypukla wady rozwiązań emulowanych, ale ogranicza ocenę samej czystej wydajności warstwy dyskowo-sieciowej tych baz.
\end{itemize}
